\cotefile{77}
\scanpage{144}
(projectif lisse) $0$-connexe (Cela fournit la factorisation de Stein de $\underline{G}_H\longrightarrow S$).

Plus généralement, si $H'\subset H$, $d'<d$, le morphisme \deleted{\unsure{prend}} correspondant
\[
\underline{G}_H\longrightarrow\underline{G}_{H'}\times_S\underline{\mathrm{Arf}}\tag*{(2)}
\]
est (projectif lisse) $0$-connexe (\unsure{cela|ceci|donc} fournit la factorisation de Stein de $\rho_{H',H}$). cf. (1)

d) \deleted{\unsure{Soit}} Supposons tjrs $N=2n$, le morphisme
\[
\underline{G}_{\{n,n-1\}}\longrightarrow\underline{G}_{\{n-1\}}\tag*{(3)}
\]
est étale fini de degré $2$, plus précisément \deleted{la factorisation de} morphisme (2) dans ce cas définit un iso
\[
\underline{G}_{\{n,n-1\}}\xrightarrow{\sim}\underline{G}_{\{n-1\}}\times\underline{\mathrm{Arf}}
\]
(pour $N=2n$)

e) \deleted{Soit} \deleted{\unsure{K}=}%
\[
\underline{\mathrm{Arf}}\,\unsure{\widetilde S'|S'|{}'}=\underline{M}((-1)^n\delta(q)^{-1})\tag*{(4)}
\]
\deleted{est} le sous-schéma de $\underline{\det} E$ défini par $\omega^{\otimes 2}=(-1)^n\delta(q)^{-1}$, Il \deleted{trouve} \inserted{existe} un morphisme unique
\[
\underline{\mathrm{Arf}}\longrightarrow\underline{\mathrm{Arf}}'\tag*{(5)}
\]
tel que le composé avec $\underline{G}_{\{n\}}\longrightarrow\underline{\mathrm{Arf}}$ soit
\[
\left\{\begin{aligned}
&\underline{G}_{\{n\}}\longrightarrow\mathrm{Arf}'\\
&F\longmapsto\omega_F\in\underline{\det}(E)^*\simeq\mathrm{Isom}(\underbrace{\det F\det\check F}_{\underline{\mathcal O}_S},\underline{\det} E)
\end{aligned}\right.\tag*{(6)}
\]
% archive sheet number: 143
\scanpage{145}
déduit de la structure d'extension. \deleted{Cet homoph.} (5) est compatible avec les op. du groupe $\mathbb Z/2\mathbb Z$, et définit donc un iso
\[
\underline{\mathrm{Arf}}'\simeq\underline{\mathrm{Arf}}\mathbin{\wedge}_{(\mathbb Z/2)_S}\mu_{2S}\tag*{(7)}
\]
--- donc, si $2$ inv. sur $S$, un iso.
\[
\underline{\mathrm{Arf}}\xrightarrow{\sim}\underline{\mathrm{Arf}}'=\underline M((-1)^n\delta(q^{*})^{-1})
\]

\deleted{f) Les $\underline{G}_H$ sont des espaces homogènes sous $\mathrm O(q)$}

\begin{center}\rule{.24\linewidth}{.4pt}\\[-.5ex]\rule{.15\linewidth}{.4pt}\end{center}

\uline{Application aux groupes orthogonaux}

\uline{Thm} a) Supposons $N=2n+1$, Alors
\[
\underline{\mathrm O}(q)\simeq\mu_2\times\underline{\mathrm{SO}}(q)
\]
et $\underline{\mathrm{SO}}(q)$ est \inserted{(affine)} lisse à fibres conn.

\deleted{b) $\underline{\mathrm{SO}}(q)$ \unsure{b)} $\mathrm{O}(q)$} opère sur chaque $\underline{G}_H(E,q)$ \inserted{via $\underline{\mathrm{SO}}(q)\simeq\underline{\mathrm O}(q)/\mu_2$,} \deleted{et le stabilis.} en esp. hom. \inserted{\unsure{sur|sous|$\to$} $\underline{\mathrm{SO}}(q)$, le stabilisateur d'une section est un ss-groupe fermé lisse à fibres conn.}%
% slanted

La décomposition en produit est vue avec l'équat. du cor. par ailleurs, notons seulement \deleted{\unsure{qu'on}} ,

\uline{Cor} La projection $\underline{\mathrm O}(q)\longrightarrow\mu_2$ est induite par dét. : $\underline{\mathrm O}(q)\longrightarrow\mathrm G_m$.

\deleted{Pour vérifier la lissité et la connexité de $\underline{\mathrm{SO}}$, \inserted{on est ramené à b).} \uline{\unsure{Obs|$\mathrm O(q)$}} qu'on a un $F\in G_{\{n\}}(E,q)$.}
% archive sheet number: 144
\scanpage{146}
\deleted{Alors $\mathrm{SO}(q)$ est fibré sur}

\deleted{$[$Considérons}
\[
\begin{aligned}
&\deleted{u\longmapsto u(F)}\\
&\deleted{\mathrm{SO}(q)\longrightarrow\underline G_{\{n\}}}
\end{aligned}
\]
\deleted{il suffit de prouver que ce morphisme est lisse à fibres géom $0$-conn. Plus gén., On va prouver que pour tt}
\deleted{\uline{Lissité}}
\[
\deleted{D\in\underline G_H(E)}\qquad
\left\{\begin{aligned}
&\deleted{\mathrm{SO}(q)\longrightarrow\underline G_H}\\
&\deleted{u\longmapsto u(D)}
\end{aligned}\right.
\]
\deleted{est lisse à fibres géom $0$-conn. (donc $\underline G_H$ espace homogène sous $\mathrm{SO}(q)$). Pour transitivité, \unsure{SPS|OPS} $H=[1,n]$}

\deleted{Comme $\mu_2$ opère \unsure{trivialmt|transitivement} sur les $\underline G_H$, $\underline{\mathrm O}(q)$ opère via $\underline{\mathrm{SO}}(q)$}
\deleted{\uline{Lissité} \unsure{groupe des translations}. Posons}
\[
\deleted{D=\{E_1\subset\ldots\subset E_n\subset\underbrace{E}_{\unsure{S}}\}}\qquad\deleted{\mathrm{rg}\,E_i=1.}
\]
\deleted{\unsure{Suggestion|Suggested}} \deleted{\unsure{q}} \deleted{Soit}

\deleted{On voit aisément \inserted{(cf. lemme plus bas)} qu'il est associé à une base standard, $e_0,e_1\ldots e_n,e_1^*\ldots e_n^*,$ \deleted{$e_0$} avec $E_i$ engendré par $e_1\ldots e_i$. \unsure{Dé|A}}

\deleted{\unsure{À|Si} un $D'=(E'_1\subset\ldots\subset E'_n)\in G_{[1,n]}$ est associé : $e'_0,e'_1\ldots e'_n\ e_1^{\prime *}\ldots e_n^{\prime *}$. Il y a un autom. (unique) \inserted{de} $(E,q)$ $e_i\longrightarrow e'_i$, $e_i^*\longrightarrow e_i^{\prime *}$, Il donne}
% archive sheet number: 145
\scanpage{147}
\uline{Proposition} $(E,q)$ forme quadratique régulière sur schéma $S$, \inserted{$\varphi$ forme bil. associée,} $F$ sous-\deleted{fibré} \inserted{Module strict} tot. \uline{isotrope} de $E$.

a) L'hom. composé
\[
E\xrightarrow[\substack{\text{défini}\\\text{par }\varphi}]{\Psi}E^\vee\xrightarrow{\mathrm{restr}}F^\vee
\]
est épi, donc son noyau $F^\perp$ est un sous-\deleted{fibré} Module strict de $E$, dont la formation commute à tout chgt de base.

b) \inserted{Soit $G=E/F^\perp$.} La forme $\varphi$ induit un accouplement
\[
F\times G\longrightarrow\underline{\mathcal O}_S
\]
qui est une dualité, et induit des isom.
\[
G\xrightarrow{\sim}\check F,\quad F\xrightarrow{\sim}\check G
\]

\deleted{c) Soit $\underline{\mathrm{Split}}(E,F^\perp)$ \inserted{ou simplement $\underline{\mathrm{Split}}$,} le faisceau des scindages de la suite exacte}
\diagram{\begin{tikzcd}[ampersand replacement=\&,row sep=tiny,column sep=small]
\deleted{0}\arrow[r]\&\deleted{F^\perp}\arrow[r,dash]\&\deleted{E}\arrow[r]\&\deleted{G}\arrow[r]\&\deleted{0},\\
\&\deleted{\cup}\\[-.8em]
\&\deleted{F}
\end{tikzcd}}
\deleted{qui est un torseur sous $\underline{\mathrm{Hom}}(G,F^\perp)$, soit $\underline{\mathrm{Split}}\,\underline{\mathrm{isot}}(E,F^\perp)$ \inserted{ou $\underline{\mathrm{Splitisot}}$,} le sous-faisceau des scindages \deleted{\unsure{$\sigma$}} \uline{isotropes}, i.e. tels que $\theta(G)$ soit \deleted{tot.} isotrope dans $E$, \deleted{\unsure{Cons}}}
\[
\deleted{\underline{\mathrm{Splitisot}}(E,F^\perp)\lhook\joinrel\longrightarrow\underline{\mathrm{Split}}(E,F^\perp)}.
\]
\deleted{Considérons d'autre part la suite exacte}
\[
\deleted{0\longrightarrow\underline{\mathrm{Hom}}(G,F)\longrightarrow\underline{\mathrm{Hom}}(G,F^\perp)\longrightarrow\underline{\mathrm{Hom}}(G,Q)\longrightarrow0}
\]
\deleted{où}
\[
\deleted{Q=F^\perp/F}.
\]
% archive sheet number: 146
\scanpage{148}
c) Soit $Q=F^\perp/F$, alors $q$ induit par passage au quotient une forme \deleted{\unsure{Q}} $q_Q$ sur $Q$, et \deleted{la forme bil. associée est} \inserted{celle-ci est régulière}

d) Soit
\[
\theta:G\longrightarrow E
\]
un scindage de la suite exacte
\[
0\mathbin{-}F^\perp\longrightarrow E\longrightarrow G\longrightarrow0\tag*{($*$)}
\]
\deleted{\unsure{et}} et considérons la forme quadratique
\[
q_\theta:F\times G\longrightarrow\underline{\mathcal O}_S
\]
\unsure{définie|défini} par
\[
q_\theta(x,y)=q(x+\deleted{\unsure{\varphi}}\theta(y))=q(\theta(y))+\deleted{\unsure{\varphi}}\langle y,x\rangle
\]
La forme bilinéaire associée est
\[
\begin{aligned}
\varphi_\theta(x,y\,;x',y')&=\varphi(x+\theta y,\theta y')+\varphi(x',\theta y)\\
&\inserted{=\varphi(\theta y,\theta y')+\langle y',x\rangle+\langle y,x'\rangle}
\end{aligned}
\]
et elle est non \uline{dégénérée}, donc $q_\theta$ est une forme quadratique régulière (de rg pair). Si $Q_\theta$ est l'orthogonal de $F+\theta(G)$, on a \deleted{donc}
\[
E\simeq F\oplus Q_\theta+\theta(G)
\]
\inserted{et \deleted{\unsure{$Q_\theta$ forme|la forme}} \inserted{la forme} $q|Q_\theta$ est régulière. $[$\uline{NB} $Q_\theta\simeq Q=F^\perp/F$, et la forme $q_Q$ déduite de $q_\theta$ ne dépend pas du choix de $\theta$, étant obtenue par passage au quotient\ldots}%
% slanted

d) Supposons $S$ affine, \deleted{\unsure{s}} ce qui assure l'existence d'un $\theta\ldots$ \deleted{\unsure{s}}. Alors il existe (pour $\theta$ fixé)
\[
u:G\longrightarrow F
\]
tel que, posant
\[
\theta_u^{\deleted{\unsure{s}}}=\theta+u
\]
on ait
\[
q\circ\theta_u=0
\]
\margin{\uline{NB} condition qui équivaut à $(q\circ\theta)(y)+\langle y,u(y)\rangle=0$ i.e. $B_u(y,y)=-(q\circ\theta)(y)$ où $B_u$ est la forme bil. associée à $u$}
i.e. $\theta_u(G)$ soit tot. isotrope dans $E$ (scindage isotrope). L'ensemble des \deleted{\unsure{sections (?)|$u\in\mathrm{Hom}($}}
\[
u\in\mathrm{Hom}(G,\underset{\substack{\shortparallel\\\check G}}{F})\simeq\mathrm{Bil}(G,G)
\]
tels que \deleted{$q\circ$} $\theta_u$ soit isotrope, est un torseur sous le sous-groupe $\mathrm{Alt}(G,G)$ de $\mathrm{Bil}(G,G)$.

Ainsi (indépendamment du choix d'un $\theta$)
% archive sheet number: 147
\scanpage{149}
l'ensemble $\mathrm{Scind}\,\mathrm{isot}(E,F^\perp)$ des scindages isotropes de $(*)$ est stable \unsure{pour|par} \deleted{\unsure{B}} les opérations de $\mathrm{Alt}(G,G)\hookrightarrow\mathrm{Bil}(G,G)\simeq\mathrm{Hom}(G,F)\subset\mathrm{Hom}(G,F^\perp)$, et le quotient $\unsure{\mathrm{Sind}\,|\,\mathrm{Scind}\,}\mathrm{isot}/\mathrm{Alt}(G,G)$ est isomorphe (par la flèche canonique) à $\mathrm{Scind}/\underbrace{\mathrm{Bil}(G,G)}_{\substack{\shortparallel\\\mathrm{Hom}(G,F)}}$ (qui est lui-$=$ un torseur sous $\mathrm{Hom}(G,Q)$, où $Q=F^\perp/F$\ldots

\deleted{\unsure{d)}} e) Soit $\theta$ un scindage isotrope de $(*)$. Alors $(E,q)$ est somme directe quadratique de $F+\underbrace{\theta(G)}_{\substack{\shortparallel\\\check F}}$, \deleted{\unsure{avec sa}} \deleted{décomposé en \unsure{form}} \inserted{où $q$ est} sous forme décomposée $q(x,x')=\langle x,x'\rangle$, et de $Q_\theta$.

\deleted{\uline{Corollaire} Considérons un drapeau}
\[
\deleted{D:\quad0\subsetneq E_1\subsetneq E_2\subset\ldots\subsetneq E_{\unsure{r|\nu}}=F\qquad E.}
\]
\deleted{d'où un drapeau complété}
\[
\deleted{0\subsetneq E_1\ldots\subsetneq E_r=F\subset F^\perp\subsetneq E_{r-1}^{\perp}\subsetneq\ldots E_1^\perp\subsetneq F}
\]
\deleted{[\uline{NB} on peut avoir $F=F^\perp$ donc $Q=F^\perp/F=0$ \ldots]. Un scindage isotrope de $(*)$ définit alors une structure suppl. dans $E$ identifié à $F\oplus\check F\oplus Q$}
% archive sheet number: 148
\scanpage{150}
\uline{Corollaire} 1 \unsure{Soient} $(E,q),(E',q')$ formes quad. rég. de $\overline{\mathrm m}$ rg \inserted{$N=2n$ ou $N=2n+1$} sur $S$, \deleted{$F\unsure{\in|\subset}E$} et $H\subset[1,n]$, enfin $D\in G_H(E,q)$, $D'\in G_H(E',q')$. \deleted{Alors si} $D=\{E_1,\ldots,E_{\unsure{r|\nu}}\}$, $D'=\{E'_1,\ldots,E'_r\}$. \deleted{Supposons $S$ affine, et supposons} Posons $Q=E_r^\perp/E_r$, $Q'=E_r^{\prime\perp}/E'_r$, de sorte que $q_\theta$ et $q_{\theta'}$ sont des formes quad. de $=$ rang sur $S$. Si $q_\theta\simeq q_{\theta'}$, \inserted{$\unsure{R|P}_i=E_i/E_{i-1}\simeq R'_i=E'_i/E'_{i-1}$,} et si $S$ est affine, alors \deleted{$q\unsure{\simeq}q'$} $\exists\,u:(E,q)\simeq(E',q')$ tel que $u(D)=D'$

On en conclut

\uline{Corollaire} 2 Il existe loc (ét) un isom $u:(E,q)\longrightarrow(E',q')$ tel que $u(D)=D'$.

Donc,

\uline{Corollaire 3} $\underline G_H(E,q)$ est un espace homogène (ét) sous $\mathrm{Aut}(E,q)=\mathrm O(q)$. En d'autres \uline{termes}, si $D\in G_H(E,q)$, alors
\[
\begin{aligned}
&\mathrm O(q)\longrightarrow\underline G_H\\
&u\longmapsto u(D)
\end{aligned}
\]
est un épim (ét), \deleted{identifiant} définissant un isom
\[
\underline G_H\xleftarrow{\sim}\mathrm O(q)/\unsure{\widetilde{\mathrm O}|\mathrm B}_{D}(q)
\]
où $\widetilde{\mathrm O}_{D}(q)$ est le stabilisateur de $D$ dans $\mathrm O(q)$.
% archive sheet number: 149
\scanpage{151}
\deleted{\uline{Corollaire}}

Nous pouvons préciser le cor. 1 ainsi :

\uline{Cor 4} Sous les conditions du cor. 1, posons \inserted{pour $1\leq i\leq\unsure{r|\nu}$} \deleted{donnons nous pour tt $1\leq i\leq\unsure{r|h|\nu}$ un iso}
\[
\unsure{R|P}_i=E_i/E_{i-1},\qquad R'_i=E'_i/E'_{i-1}\qquad(E_0=0,E'_0=0)
\]
et donnons nous des isom
\[
\begin{aligned}
u_i:\ &R_i\xrightarrow{\sim}R'_i\\
u_q:\ &Q\xrightarrow{\sim}Q'
\end{aligned}
\]
avec $v$ \unsure{orthogonale|orthogonal} pour $q,q'$. Alors, \deleted{il existe} \inserted{$S$ étant affine,} \inserted{a) il existe} un isom
\[
u:(E,q)\longrightarrow(E',q')
\]
tel que $u(D)=D'$, et que $u$ induise les $u_i$ et $u_q$.

\deleted{\unsure{b)}} L'ens de ces $u$ est \inserted{(\unsure{évidemt|évident})} un torseur à dr. sous le \inserted{ss} groupe \inserted{$\underline{\mathrm O}_D(q)^+$,} des \deleted{\unsure{auto|$\to$} $\mathrm O_D(q)(S$} $\mathrm{Aut}(E,q,D)=\underline{\mathrm O}_D(q)(S)$, formé des autom qui induisent l'identité sur les $G_i$ et sur $Q$.

\deleted{Considérons d'autre} \inserted{Considérons d'autre} part l'ens \deleted{des} produit
\[
\Sigma=\prod_{1\leq i\leq r}\Sigma_i\times\Sigma_0
\]
où $\Sigma_i$ est l'ens des scindages de
\[
0\longrightarrow E_{i-1}\mathbin{-}E_{\deleted{\unsure{r}}i}\mathbin{-}R_i\mathbin{-}0
\]
et $\Sigma_0$ l'ens des scindages \uline{isotropes} de
\[
0\longrightarrow\deleted{\unsure{E}}E_r\longrightarrow E_{r\deleted{\unsure{2}}}^{\perp}\longrightarrow Q\longrightarrow0
\]
\deleted{\unsure{On}} Considérons \deleted{d'autre part} \inserted{aussi}
\[
\overline E=\bigoplus(R_i+\check R_i)\oplus Q
\]
avec la forme $\overline q$ somme des formes
% archive sheet number: 150
\scanpage{152}
évidentes $\langle x_i,x'_i\rangle$ sur les $G_i\oplus\check G_i$, et $q_Q$ sur $Q$. \deleted{Alors} Soit $\overline D=\{\overline E_1,\ldots,\overline E_{\unsure{r|\nu}}\}$, avec $\overline E_i=\sum_{1\leq j\leq i}\unsure{R|P}_j\in\overline E$. \deleted{Considérons}

Considérons
\[
\left\{\begin{aligned}
\overline R_i&=\overline E_i/\overline E_{i-1}\\
\overline Q&=\overline E_r^\perp/\overline E_r
\end{aligned}\right.
\]
on a donc des isom canoniques
\[
\begin{aligned}
\overline R_i&\xrightarrow{\overset{\bar u_i}{\sim}}R_i\\
\overline Q&\mathrel{\overset{\unsure{\bar v|\bar u|\bar u_q}}{\simeq}}Q,
\end{aligned}
\]
le dernier étant compatible avec $\bar q,q$.

Ceci posé, on a

\uline{Cor. 5} L'ens $\Sigma$ est canoniquement isom à l'ens. des isom $(\overline E,\bar q,\overline D)\longrightarrow(\overset{\deleted{\unsure{-}}}{E},q,D)$ compatibles avec les $\bar u_i,\bar v$.

En termes schématiques, le schéma produit $\underline\Sigma=\prod_{1\leq\unsure{j|i}\leq r}\underline\Sigma_i\times\underline\Sigma_0$ est canoniquement isomorphe au schéma
\[
\underline{\mathrm{Isom}}_{\{\bar u_i\},\bar v}\big((\overline E,\bar q,\overline D);(E,q,D)\big)
\]
\margin{sans hyp. $S$ affine !}%
% slanted
Ce dernier est un torseur à gauche sous $\underline{\mathrm O}_D^+(q)$ (à droite sous $\underline{\mathrm O}_{\overline D}^{+}(\bar q)$)

Comme $\underline\Sigma$ est \unsure{évidemt|évidemment} lisse \inserted{$0$-}connexe sur $S$, on en conclut
% archive sheet number: 151
\scanpage{153}
\uline{Cor.} 6 Le \inserted{schéma en} groupes $\underline{\mathrm O}_D(q)^+$ est lisse à fibres connexes, plus précisément, \deleted{\unsure{pour|si}} \inserted{qd $S$ est affine} (donc $\underline\Sigma$ a une section donc $\underline{\mathrm O}_D(q)^+\simeq\underline\Sigma$) il est isom. comme schéma (mais sans structure de groupe) à un fibré \deleted{affine} vectoriel
\[
\check{\mathbb V}\left(\sum_{1\leq i\leq\unsure{r|\nu}}\underline{\mathrm{Hom}}(\unsure{R|P}_i,E_{i-1})\oplus\sum_{1\leq i\leq r}\underline{\mathrm{Hom}}(\overline G_{\unsure{i|r}},\overline F^\perp)^{!}\right)
\]
où $\underline{\mathrm{Hom}}(\overline G,\overline F^\perp)^{!}$ est le noyau de l'épi
\[
\begin{aligned}
&\underline{\mathrm{Hom}}(\overline G,\overline F^\perp)\longrightarrow\underline{\mathrm{Quad}}(\overline G)\\
&u\longmapsto\deleted{(\unsure{\bar y|\bar q}\mapsto\bar q(\unsure{\bar u}(\bar y))+\varphi(\unsure{-})}\\
&\hphantom{u\longmapsto}(\bar y\longmapsto\bar q(\bar y+u(\bar y)))
\end{aligned}
\]

Calculons les dimensions, posant \deleted{dim}
\[
\begin{aligned}
&\mathrm{rg}\,R_i=d_i,\quad \mathrm{rg}\,Q=\deleted{\unsure{d}}\rho
\qquad\text{donc }\underline{N=2\overbrace{\sum\deleted{\unsure{i}}d_i}^{d}+\rho}\\
&\mathrm{rg}\,\overline F=\mathrm{rg}\,G=\underbrace{\sum d_i}_{\deleted{d}}\overset{\mathrm{dfn}}{=}d
\end{aligned}
\]
\[
\left\{\begin{aligned}
\dim_S\underline\Sigma_i&=d_i\sum_{1\leq j\leq i-1}d_j\\
\dim_S\underline\Sigma_0&=\dim_S\underline{\unsure{\mathrm{Morph}|\mathrm{Hom}}}(\overline G,\overline F^\perp)-\dim_S\underline{\mathrm{Quad}}(\overline G)\\
&=\deleted{\unsure{\sum d_i}}d(d+\rho)-\frac{d(d+1)}{2}=\frac{d(d-1)}{2}+d\rho
\end{aligned}\right.
\]
\[
\begin{aligned}
\dim_S\underline{\mathrm O}_D(q)^+
 &=\underbrace{\sum_{1\leq j<i\leq r}d_i d_j}_{\deleted{\unsure{=}}\ \frac12(d^2-\sum d_i^2)}+\frac{d(d-1)}2+d\rho\\
 &=d^2+d\rho-\deleted{\unsure{\tfrac12}}\sum\deleted{\unsure{i}}\frac{d_i(\deleted{\unsure{+}}d_i^{\deleted{\unsure{2}}}+1)}2
\end{aligned}
\]
\[
\dim_S\left(\prod\mathrm{Gl}(R_i)\mathbin{\times}\mathrm O(q_Q)\right)=\sum d_i^2+\frac{\rho(\rho-1)}2
\]
\[
\dim_S\underline{\mathrm O}_D(q)=d^2+d\rho+\frac{\rho(\rho-1)}2+\unsure{\tfrac12|\deleted{\tfrac12}}\sum_i\frac{d_i(\deleted{\unsure{d}}d_i-1)}2
\]
\[
\dim_S\underline G_D(q)=
\underbrace{\dim_S\underline{\mathrm O}_{\deleted{\unsure{D}}}(q)}_{\frac{N(N-1)}2=2d^2+\deleted{\unsure{2}}d(2\rho-1)+\frac{\rho(\rho-1)}2}
-\dim_S\underline{\mathrm O}_D(q)=\frac{\deleted{\unsure{N(N-1)}}}{2}
\]
% archive sheet number: 152
\scanpage{154}
\[
\dim_S\underline G_H(q)=d^2+d\rho-d-\sum\frac{d_i(d_i-1)}2
\]

Par exemple
\[
\dim_S\underline G_{\{d\}}(q)=\frac{d(d+1)}{2}+d(\rho\deleted{\unsure{+}}-1)\qquad(\rho=N-2d)
\]
\deleted{\unsure{Ai|Ainsi}}Ainsi, si $\unsure{q|f}$ est de rg $N=2n$, et $d=\{n\}$, \inserted{(donc $\rho=0$)} on trouve
\[
\frac{\deleted{\unsure{d}}n(n+1)}{2}\deleted{\unsure{+|-}}-n=\frac{n(n-1)}{2}
\]

\begin{center}\rule{.14\linewidth}{.4pt}\end{center}

Supposons $N=2\underset{\deleted{\unsure{0}}}{n}+1$ --- de sorte que
\[
\mathrm O(q)=\mathrm{SO}(q)\times\mu_2,
\]
\deleted{On va vérifier} \inserted{où}
\[
\left\{\begin{aligned}
&\mu_2=\mathrm O(q)\cap G_m=\mathrm Z(\mathrm O(q))\\
&\underline{\mathrm{SO}}(q)=\underline{\mathrm O}(q)\cap\underline{\mathrm{Sl}}(E)
\end{aligned}\right.
\]
\deleted{\unsure{$\Rightarrow$ il résulte}} \deleted{que $\mathrm{SO}(q)$ est} $\mathrm O(q)$ opère sur les $\underline G_H(q)$ avec $\mu_2$ opérant trivialement, donc il opère via le quotient $\mathrm{SO}(q)$ --- donc les $\underline G_H$ sont aussi des espaces homogènes sous $\mathrm{SO}(q)$. \deleted{\unsure{Pour le|Pour}} \deleted{\unsure{dès que}} \inserted{Pour} $\forall\,D\in G_H$, \inserted{soit $\underline{\mathrm{SO}}_D(q)$} le stabilisateur dans $\underline{\mathrm{SO}}(q)$, \deleted{donc}. On a donc
\[
\underline{\mathrm{SO}}_D(q)=\underline{\mathrm O}_D(q)\cap\underline{\mathrm{Sl}}(E)
\]
Considérons l'hom canonique
\[
\underline{\mathrm O}_D(q)\xrightarrow{\alpha}\deleted{\prod_{1\leq i\leq\unsure{n|r|\nu}}\underline{\mathrm{Gl}}(\unsure{R|P}_i)}\times\underline{\mathrm O}(q_Q)
\]
(où, comme dessus, $\deleted{\unsure{q_Q}}Q=E_{\unsure{r|\nu}}^\perp/E_r$, $D=\{E_1\subset\ldots\subset E_r\}$). On constate que
% archive sheet number: 153
\scanpage{155}
l'on a commutativité dans
\diagram{\begin{tikzcd}[ampersand replacement=\&,column sep=large]
\underline{\mathrm O}_D(q)\arrow[r]\arrow[rr,bend right=24,"\det_E"']\&\mathrm O(q_Q)\arrow[r,"\det_Q"]\&G_m
\end{tikzcd}}
donc
\[
\underline{\mathrm{SO}}_D(q)=\alpha^{-1}(\mathrm{SO}(q_Q))
\]

On trouve donc

\uline{Proposition} \deleted{$\underline{\mathrm{SO}}_D(q)$ est \unsure{le noyau|isom.}} \inserted{On a en tous cas} \unsure{des suites|la suite|On a} \deleted{\unsure{exacte|$\alpha$}} \deleted{$\underline{\mathrm{SO}}_D(q)\longrightarrow\prod_{1\leq i\leq\unsure{r|\nu}}\mathrm{Gl}(P_i)\times\underline{\mathrm{SO}}_{\deleted{\unsure{D}}}(q_Q)$} \deleted{induit un}
\diagram{\begin{tikzcd}[ampersand replacement=\&,column sep=small,row sep=large]
1\arrow[r]\&\underline{\mathrm O}_D(q)^+\arrow[r]\&\underline{\mathrm O}_D(q)\arrow[r]\&\prod_{1\leq i\leq r}\mathrm{Gl}(P_i)\times\mathrm O(q_Q)\arrow[r]\&1\\
1\arrow[r]\&\underline{\mathrm{SO}}_D(q)^+\arrow[u,"\sim"']\arrow[r]\&\underline{\mathrm{SO}}_D(q)\arrow[u,hook]\arrow[r]\&\prod_{1\leq i\leq r}\mathrm{Gl}(P_i)\times\mathrm{SO}(q_Q)\arrow[u,hook]\arrow[r]\&1
\end{tikzcd}}
donc $\underline{\mathrm{SO}}_D(q)$ est une extension de $\prod\mathrm{Gl}(P_i)\times\underline{\mathrm{SO}}(q_Q)$ par $\underline{\mathrm{SO}}_D(q)^+\simeq\underline{\mathrm O}_D(q)^+$, \deleted{\unsure{donc|d'où}} \inserted{qui} est \inserted{lisse à fibres \unsure{géom. conn.|connexes}} \deleted{est lisse et à fibres connexes (comme ext. de $2$ gr qui ont ces propriétés}.

\uline{Thm} \inserted{Supposons $E$ de rg impair $N=2n\unsure{.}$, Alors} $\underline{\mathrm{SO}}(q)$ $(=\underline{\mathrm O}(q)\cap\underline{\mathrm{Sl}}(E))$ est un groupe lisse à fibres connexes.

\uline{OPS} par descente qu'il existe $F\subset E$ isotrope de rg $n$, donc $\deleted{\unsure{E}}F^\perp/F=Q$ de rg $1$, donc $\mathrm{SO}(q_Q)=1$, donc \deleted{\unsure{on}} (ici $D=\{F\}$) $\mathrm{SO}_D(q)$ est extension de $\deleted{\unsure{\prod_i}}\mathrm{Gl}(F)$ (qui est lisse à fibres conn) par $\underline{\mathrm{SO}}_D$ (qui est itou)
% archive sheet number: 154
\scanpage{156}
donc est lisse à fibres connexes. D'autre part $\underline{\mathrm{SO}}(q)/\underline{\mathrm{SO}}_D(q)\simeq\underline G_{\{n\}}(q)$ est lisse à fibres géom connexes, donc aussi $\underline{\mathrm{SO}}(q)$.

\uline{Corollaire} Les $\underline{\mathrm{SO}}_D(q)$ \inserted{$=B_D(q)$} sont lisses à fibres connexes. On a $\underline{\mathrm O}_D(q)\simeq B_D(q)\times\mu_2$.

\begin{center}\rule{.12\linewidth}{.4pt}\end{center}

Supposons enfin $N=2n$, de sorte que $\mu_2=\mathrm Z(\mathrm O(q))\subset\mathrm{Sl}(E)$. On ne va pas définir $\underline{\mathrm{SO}}(q)$ comme $\mathrm O(q)\cap\mathrm{Sl}(E)$, à cause de car. $2$. Mais on fait opérer par transp de str. $\underline{\mathrm O}(q)$ sur $\underline{\mathrm{Arf}}(q)$, donc on a un hom can.
\[
\underline{\mathrm O}(q)\xrightarrow{\varepsilon_q}(\mathbb Z/2)_S
\]
\margin{explicite : soit $\eta:(\mathbb Z/2)_S\longrightarrow\mu_{2S}$ can, alors $\eta\circ\varepsilon_q=\det_q$ (restriction de $\det_E$ à $\underline{\mathrm O}(q)$) \uline{Cor} si $2$ inv. sur $S$ $\underline{\mathrm{SO}}(q)=\underline{\mathrm O}(q)\cap\mathrm{Sl}(E)$}
et on pose
\[
\underline{\mathrm{SO}}(q)=\mathrm{Ker}\,\varepsilon
\]
de sorte qu'on a la suite exacte
\[
1\mathbin{-}\underline{\mathrm{SO}}(q)\mathbin{-}\underline{\mathrm O}(q)\xrightarrow{\varepsilon_q}(\mathbb Z/2)_S\longrightarrow1
\]

\inserted{\uline{Prop} \deleted{$\mathrm{SO}(q)$} Soit $H\subset[1,n]$, \inserted{$d=\sup H$,} alors $\underline{\mathrm{SO}}(q)$ est transitif sur $\underline G_H(q)$ \deleted{\unsure{ssi|sauf si} $d\unsure{=}\sup H$} $d\ne n$, et \unsure{inv.|aussi} si $S\ne\varnothing$. Si $d=n$, et si on a une section de $\underline{\mathrm{Arf}}(S)$ sur $S$, alors $\underline{\mathrm{SO}}(q)$ est transitif sur $\underline G_H^\omega(q)$.}%
% margin note written sideways

Considérons comme dessus un $D$, d'où
% archive sheet number: 155
\scanpage{157}
$\underline{\mathrm O}_D(q)$ et
\[
1\longrightarrow\underline{\mathrm O}_D(q)^+\longrightarrow\underline{\mathrm O}_D(q)\longrightarrow\prod_{1\leq i\leq\unsure{r|\nu}}\mathrm{Gl}(P_i)\times\mathrm O(q_Q)\mathbin{-}1
\]
On vérifie alors

\deleted{\uline{Proposition} Posons $\underline{\mathrm{SO}}_D(q)=\underline{\mathrm O}_D(q)\cap\underline{\mathrm{SO}}$, stabilisateur de $D$ dans $\underline{\mathrm{SO}}(q)$. Alors \unsure{S}}

\uline{Prop} La restriction $\varepsilon_q|\underline{\mathrm O}_D(q)$ est égal au composé $\underline{\mathrm O}_D(q)\longrightarrow\underline{\mathrm O}(q_Q)\xrightarrow{\varepsilon_{q_Q}}(\mathbb Z/2)_S$, donc $\underline{\mathrm{SO}}_D(q)$ est l'image inverse de $\underline{\mathrm{SO}}(q_Q)$ par $\underline{\mathrm O}_D(q)\longrightarrow\underline{\mathrm O}(q_Q)$, et on a un hom de suites exactes

\deleted{\uline{Cor} $\underline{\mathrm{SO}}_D(q)$ est une extension de \unsure{$\prod\mathrm{Gl}$}}
\diagram{\begin{tikzcd}[ampersand replacement=\&,column sep=small,row sep=large]
1\arrow[r]\&\underline{\mathrm O}_D(q)^+\arrow[r,dash]\arrow[d,"\sim"]\&\underline{\mathrm O}_D(q)\arrow[r]\&\prod_{1\leq i\leq r}\mathrm{Gl}(P_i)\times\underline{\mathrm O}(q_Q)\arrow[r]\&1\\
1\arrow[r]\&\underline{\mathrm{SO}}_D(q)^+\arrow[r,dash]\&\underline{\mathrm{SO}}_D(q)\arrow[u,hook]\arrow[r]\&\prod_{1\leq i\leq r}\mathrm{Gl}(P_i)\times\underline{\mathrm{SO}}(q_Q)\arrow[r,dash]\&1,
\end{tikzcd}}
de sorte que $\underline{\mathrm{SO}}_D(q)$ est une ext. de $(\prod\mathrm{Gl}(P_i))\times\mathrm{SO}(q_Q)$ par $\underline{\mathrm{SO}}_D(q)^+\simeq\underline{\mathrm O}_D(q)^+$, qui est lisse à fibres conn.

\uline{Thm} $(E,q)$ forme quad. rég. de rg pair $2n$. Alors $\underline{\mathrm S\deleted{\unsure{\mathrm O}}\unsure{\mathrm E|\mathrm O}}(q)$ est lisse à fibres connexes.

OPS $\exists\,F\subset G_n(q)$, d'où section $\omega$ de $\underline{\mathrm{Arf}}(q)$ sur $S$. Prenons ci-dessus $D=\{F\}$, \deleted{$F$ est de rg $n$,} donc $Q=0$. Donc par ce qui précède, $\underline{\mathrm{SO}}_D(q)$ est lisse à fibres connexes.
% archive sheet number: 156
\scanpage{158}
D'autre part $\underline{\mathrm{SO}}(q)/\underline{\mathrm{SO}}_D(q)\simeq\underline G_{\{n\}}^\omega$ est lisse à fibres connexes, donc aussi $\underline{\mathrm{SO}}(q)$, cqfd.

\uline{Corollaire} Les $\underline{\mathrm{SO}}_D(q)=B_D(q)$ sont lisses à fibres connexes.
% archive sheet number: 157
